% abstract em inglês
\begin{resumo}[Abstract]
 \begin{otherlanguage*}{english}
   This work presents a comparative and diagnostic analysis of the performance of heterogeneous enterprise Wi-Fi networks using physical and logical telemetry. 
   The study is based on a dataset collected continuously over 18 days in a high-density university campus, covering two distinct network architectures (Ruckus Wireless and Ubiquiti UniFi), resulting in over 350,000 processed telemetry records. 
   The proposed methodology involves a structured pipeline of data cleaning, standardization, statistical hypothesis testing, and heuristic and unsupervised classification of network bottlenecks. 
   Non-parametric Mann-Whitney tests revealed statistically significant performance disparities ($p < 0.001$) between the 2.4~GHz and 5~GHz bands, as well as severe discrepancies in frame retransmission rates between manufacturers (UniFi averaging 22.1\% vs. Ruckus averaging 5.7\%). 
   Bottleneck classification indicated that client throughput was severely degraded by link quality issues (up to 89.3\% throughput drop under weak signal) and roaming instability, which reduced Ruckus client throughput by 38.8\%. 
   Finally, the K-Means clustering algorithm ($K=4$) was validated using Silhouette (0.3383) and Davies-Bouldin (0.7953) metrics for the Ruckus network, proving the effectiveness of segmenting Access Points into distinct operational profiles and detecting anomalies in the physical transmission medium. 
   This research contributes a reproducible and scalable methodology for automated diagnosis and quality of service management in enterprise wireless networks.

   \vspace{\onelineskip}

   \noindent
   \textbf{Keywords}: Network Telemetry. Enterprise Wi-Fi. K-Means. Statistical Analysis. Bottleneck Diagnosis.
 \end{otherlanguage*}
\end{resumo}